\chapter{Discussion}
\label{ch:discussion}

This chapter interprets the results presented in Chapter~\ref{ch:results}. It first relates the findings back to the expectations established by the literature, then examines the implications for practitioners, and concludes by setting out the threats to validity and the consequent limitations of the work.

\section{Interpretation Against Prior Expectations}

The literature surveyed in Chapter~\ref{ch:literature} predicts qualitatively that segmentation should reduce lateral-movement opportunity at modest performance cost. The results obtained here are consistent with that prediction in direction and provide a quantitative anchor for its magnitude in the specific configuration tested. A 75-percentage-point reduction in containment failure (from 100\% exposure to 25\%) is a substantial change, and it matches the order-of-magnitude claims made in vendor literature \parencite{VMware}{2021}; \parencite{Cisco}{2022} and the qualitative statements in standards documents \parencite{NIST SP~800-41}{2009}; \parencite{NIST}{2024}.

The result on policy accuracy is particularly informative. \textcite{Wool}{2004} found that almost all production firewalls contain policy errors; the mechanism by which the present project's flat firewall fails its tests is a special case of this finding, namely a complete absence of internal policy. By contrast, the segmented firewall reaches 100\% accuracy because every internal flow is enumerated explicitly under a default-deny policy. This is the same property that the principle of \emph{economy of mechanism} \parencite{Saltzer and Schroeder}{1975} predicts: simple, additive policies are easier to audit and easier to keep correct than complex, subtractive ones.

The performance result deserves separate comment. The 50~$\mu$s overhead measured here is a direct measurement of one additional Linux \texttt{iptables} forwarding decision, and is consistent with measurements published in the broader networking literature \parencite{Tanenbaum and Wetherall}{2021}. It is small enough that, for any human-facing or batch-processing workload, it is operationally invisible. The argument that ``segmentation slows everything down'' --- which the author has heard informally on more than one occasion in industry contexts --- is not supported by the data. Where segmentation \emph{does} introduce meaningful overhead in production is in the human cost of designing, documenting, and maintaining the policy, not in the per-packet runtime cost.

\section{Implications for Practice}

Several practitioner-level implications follow from the results.

\paragraph{Resource-constrained organisations should prioritise segmentation.} The implementation evaluated here uses entirely free, open-source components. The control-plane cost is the engineering effort required to design and document a zone model and an associated rule set --- not licence fees, hardware acquisition, or specialist staff. Smaller organisations that lack the budget for an SDN deployment or a microsegmentation platform can nevertheless obtain the bulk of the security benefit using \texttt{iptables}-based zone segmentation, particularly in virtualised or containerised environments. This recommendation is consistent with the position of the U.K.\ National Cyber Security Centre \parencite{NCSC}{2018}; \parencite{NCSC}{2023} and the CIS Critical Security Controls v8 \parencite{CIS}{2024}.

\paragraph{Policy correctness should be measured, not assumed.} The policy-accuracy metric defined in Section~3.4.1 is straightforward to instrument and provides an objective measure of whether a firewall actually behaves as intended. Industry experience reported by \textcite{Wool}{2004} suggests that periodic measurement against a curated test set would surface configuration drift far earlier than ad-hoc auditing. The eight-test set used here is too small for a real estate, but the methodology generalises trivially.

\paragraph{Segmentation is necessary, not sufficient.} The same-zone reachability of \texttt{student2} from \texttt{student} in the segmented topology is not a flaw --- it is a deliberate design choice consistent with the principle of zone-based segmentation \parencite{Convery}{2004} --- but it is a useful reminder that segmentation alone does not protect against an attacker whose target shares the attacker's zone. In a real environment this would correspond to a malicious or compromised student attacking another student, and would be addressed not by network controls but by host-level controls such as endpoint detection and response \parencite{Stallings}{2017} and by application-level authentication. The literature on zero-trust architecture \parencite{Rose et al.}{2020} explicitly anticipates this and is the natural next step beyond zone segmentation.

\paragraph{Segmentation supports regulatory compliance.} Several regulatory regimes either require or strongly recommend segmentation. PCI DSS v4.0 \parencite{PCI Security Standards Council}{2022} permits the cardholder-data environment to be ``isolated (segmented)'' from the rest of the corporate network as a means of reducing scope. ISO/IEC 27002 control 8.22 (\textit{Segregation of networks}) explicitly requires it \parencite{ISO/IEC}{2022}. Article 32 of the General Data Protection Regulation requires ``appropriate technical and organisational measures'' to ensure the security of processing, of which segmentation is one of the most established. The empirical results obtained here strengthen the practical case for these regulatory positions.

\section{Relationship to Zero-Trust Architecture}

It is important to be precise about what segmentation \emph{is} and \emph{is not} in the context of the modern zero-trust paradigm \parencite{Rose et al.}{2020}; \parencite{Microsoft}{2023}. The segmented topology evaluated in this project is a coarse-grained, perimeter-and-zone implementation. It enforces network-layer policy at zone boundaries but does not authenticate individual flows, does not consider device posture, and does not consider user identity. In zero-trust terms it sits towards the \emph{lower} end of maturity: it removes implicit trust between zones but retains implicit trust within zones.

The methodology developed here generalises to higher-maturity implementations without alteration. The same metrics --- policy accuracy, blast-radius exposure, performance, composite reliability --- can be computed against any topology, including one in which every host is its own zone (microsegmentation \parencite{NIST}{2016}; \parencite{VMware}{2021}) or one in which authorisation is per-flow (full zero trust). What changes with maturity is the magnitude of the reduction in blast radius and the engineering complexity of maintaining the policy. The present work establishes the lower bound for the security benefit; future work could measure the marginal improvement obtained by moving to finer granularities.

\section{Threats to Validity}

The findings are subject to several threats to validity, which are stated here explicitly so that the reader can weigh them.

\paragraph{Construct validity.} The composite reliability index defined in Equation~\ref{eq:reliability} is a weighted sum, and the choice of weights (0.30, 0.35, 0.25, 0.10) is a judgement call. Different weights would produce different scores, although the qualitative ordering (segmented~$>$~flat) is robust to any non-degenerate choice of weights because the segmented topology dominates the flat topology on three of the four constituent metrics. The weights chosen are within the range recommended by \textcite{Cherdantseva and Hilton}{2013} and were fixed before the experiments were run, mitigating concerns about post-hoc tuning.

\paragraph{Internal validity.} The two topologies were constructed to be functionally equivalent except for the segmentation, but they are not literally identical. The segmented topology has additional Docker network bridges and additional firewall rules; these introduce additional forwarding hops and hence the 50~$\mu$s latency difference. This is not a confounder --- it is the cost of segmentation, properly captured --- but it should be acknowledged as a real difference in the underlying systems.

\paragraph{External validity.} The testbed is small (ten hosts) and uses synthetic services that do not replicate the application-layer behaviour of real software. The results should therefore be read as a controlled lower bound on what segmentation can achieve, not as a forecast of behaviour in any specific real network. The literature on cyber-range scaling \parencite{Yamin, Katt and Gkioulos}{2020}; \parencite{Russo, Costa and Armando}{2020} suggests that the qualitative findings should hold at larger scales, but a definitive claim would require larger experiments.

\paragraph{Threat-model validity.} The threat model in Section~\ref{sec:threat-model} is intentionally narrow. Adversaries who can compromise the firewall itself, or who can move laterally over channels other than direct IP reachability (for example, by exploiting a shared application layer such as DNS or Active Directory), are out of scope. The work does not claim that segmentation defeats every adversary --- only that it defeats the lateral-movement techniques that depend on routable network reachability, which is the majority of those catalogued by \textcite{MITRE}{2024}.

\paragraph{Performance-measurement validity.} The performance test issues sequential HTTP requests rather than concurrent load. Real production systems would face concurrent traffic that might saturate the firewall or the test target; the measurements presented here characterise per-request latency, not throughput under load. \textcite{Bonguet and Bellaiche}{2017} review the literature on denial-of-service evaluation under load; extending the harness to include concurrent and stress workloads is an obvious follow-up.

\section{Limitations of the Implementation}

Beyond the threats to validity, three concrete limitations of the implementation are worth highlighting.

First, the firewall is implemented as a single \texttt{iptables} container. In a real deployment the firewall would itself be a high-availability cluster, and the failure modes of that cluster are not represented here. The single-firewall design also means that a compromise of the gateway would defeat segmentation in both topologies, which is part of why \textcite{Anderson}{2020} recommends defence in depth rather than a single chokepoint.

Second, the firewall rules are static. Real estates change frequently, and the operational challenge of keeping the rule set in sync with the actual application architecture is a major source of policy errors in production \parencite{Wool}{2004}. The methodology supports automated re-testing, but the experimental scope does not include change-management workflows.

Third, the testbed is purely network-layer. Application-layer attacks --- SQL injection, cross-site scripting, deserialisation flaws --- proceed over permitted flows and are entirely unaffected by segmentation \parencite{Pfleeger, Pfleeger and Margulies}{2015}. This is by design: segmentation is a network-layer control, and the work measures only what is in its remit. A complete defence requires complementary controls at every layer of the stack.
